\section{Вариант задания}
\label{sec:Chapter4} \index{Chapter4}
Индивидуальные варианты заданий отличаются комбинацией граничных условий.
% Варианты приведены в следующей таблице (1Р -- условие первого рода, П -- периодическое условие) 1. 

В текущей работе рассматривается вариант 5:
\begin{itemize}
    \item $
        u_{analytical}
        =
        sin \left( \frac{2\pi}{L_x}x \right)
        \cdot
        sin \left( \frac{2\pi}{L_y}y \right)
        \cdot
        sin \left( \frac{2\pi}{L_z}z \right)
        \cdot
        cos \left( a_t \cdot t + 4\pi \right),
        a_t
        =
        \pi \sqrt{
            \dfrac{4}{L^2_x}
            +
            \dfrac{1}{L^2_y}
            +
            \dfrac{1}{L^2_z}
        }
    $,
    \item \(x\) -- периодическое условие (\ref{equation:7}),
    \item \(y\) -- условие первого рода (\ref{equation:5}),
    \item \(z\) -- условие первого рода (\ref{equation:6}).
\end{itemize}

Функция \(u_{analytical}\) напрямую зависит от введенных \(L_x\), \(L_y\), \(L_z\), которые являются входными аргументами программы.
Функция \(\varphi(x, y, z)\) зависит напрямую от \(u_{analytical}\) (\ref{equation:2}).